\renewcommand{\currentchapter}{6}
\title{Peta Materi UTBK 2027}
\subtitle{Bab 6 \textbar\ Kurikulum operasional tujuh subtes dalam StudentPro}
\date{SMAN 1 Muntilan bersama StudentPro \textbar\ Tahun Ajaran 2026/2027}

\begingroup
\setbeamercolor{background canvas}{bg=StudentNavy}
\begin{frame}[plain]
  \vspace{0.65cm}
  \kicker{SMAN 1 Muntilan \textbar\ StudentPro}
  \vspace{0.08cm}
  {\usebeamerfont{title}\color{StudentWhite}\inserttitle\par}
  \vspace{0.38cm}
  {\usebeamerfont{subtitle}\color{StudentSoft}\insertsubtitle\par}
  \vfill
  {\color{StudentGold}\rule{2.3cm}{2.5pt}}\par
  \vspace{0.20cm}
  {\usebeamerfont{author}\color{StudentWhite}\insertauthor}\par
  {\usebeamerfont{date}\color{StudentSoft}\insertdate}
  \vspace{0.42cm}
\end{frame}
\endgroup

\begin{frame}{Bab 6 mengubah framework menjadi kurikulum operasional}
  \kicker{Posisi Bab dalam Laporan}
  \begin{columns}[T,onlytextwidth]
    \begin{column}{0.30\textwidth}
      \begin{block}{Bab 2}
        Menjelaskan struktur, tujuan, jumlah soal, durasi, dan gambaran tujuh subtes.
      \end{block}
    \end{column}
    \begin{column}{0.30\textwidth}
      \begin{block}{Bab 6}
        Memetakan materi, indikator, tingkatan belajar, dan penandaan konten.
      \end{block}
    \end{column}
    \begin{column}{0.30\textwidth}
      \begin{block}{Bab 7}
        Menyajikan contoh soal, kunci, pembahasan, dan variasi bentuk soal.
      \end{block}
    \end{column}
  \end{columns}
  \vspace{0.26cm}
  \claim{Bab ini menjawab: apa yang dipelajari, dalam urutan apa, dan bagaimana kemajuannya diukur?}
\end{frame}

\begin{frame}{Kurikulum dibangun dengan empat prinsip}
  \kicker{Dasar Penyusunan Materi}
  \begin{columns}[T,onlytextwidth]
    \begin{column}{0.22\textwidth}
      \begin{block}{Selaras}
        Mengikuti kemampuan yang dinilai dalam framework resmi 2026.
      \end{block}
    \end{column}
    \begin{column}{0.22\textwidth}
      \begin{block}{Bertahap}
        Bergerak dari konsep dasar menuju penerapan dan simulasi.
      \end{block}
    \end{column}
    \begin{column}{0.22\textwidth}
      \begin{block}{Terukur}
        Setiap materi memiliki indikator dan bukti ketuntasan.
      \end{block}
    \end{column}
    \begin{column}{0.22\textwidth}
      \begin{block}{Adaptif}
        Urutan belajar menyesuaikan diagnostik dan perkembangan siswa.
      \end{block}
    \end{column}
  \end{columns}
  \vspace{0.28cm}
  \centering{\color{StudentNavy}\bfseries\fontsize{14}{17}\selectfont Materi tidak disusun sebagai tumpukan bab, tetapi sebagai jalur penguasaan.}
\end{frame}

\begin{frame}[shrink=7]{Tujuh subtes menghasilkan tujuh profil kemampuan}
  \kicker{Peta Kurikulum Utama}
  {\fontsize{9.3}{11.5}\selectfont
  \renewcommand{\arraystretch}{0.94}
  \begin{tabularx}{\textwidth}{>{\bfseries\color{StudentBlue}}p{0.24\textwidth} >{\RaggedRight\arraybackslash}X >{\RaggedRight\arraybackslash}p{0.29\textwidth}}
    \toprule
    \textbf{Subtes} & \textbf{Fokus kurikulum} & \textbf{Profil yang dihasilkan} \\
    \midrule
    Penalaran Umum & Induktif, deduktif, dan kuantitatif & Menemukan pola dan menarik simpulan \\
    PPU & Makna, hubungan gagasan, dan pengetahuan dalam konteks & Memahami bahasa dan informasi \\
    PBM & Memilih, menata, dan memperbaiki gagasan & Menyunting tulisan secara logis \\
    Pengetahuan Kuantitatif & Bilangan, aljabar, geometri, data, dan peluang & Memilih prosedur matematika efisien \\
    LBI & Menemukan, memahami, menilai, dan merefleksikan teks & Literasi menurut konteks bacaan \\
    LBE & Pemahaman teks akademik bahasa Inggris & Membaca bukti dan argumen berbahasa Inggris \\
    Penalaran Matematika & Merumuskan, menggunakan, dan menafsirkan model & Memecahkan masalah kontekstual \\
    \bottomrule
  \end{tabularx}}
\end{frame}

\begin{frame}{Setiap materi dilalui melalui lima tingkat}
  \kicker{Jalur Penguasaan}
  {\fontsize{10.0}{12.5}\selectfont
  \renewcommand{\arraystretch}{0.98}
  \begin{tabularx}{\textwidth}{>{\bfseries\color{StudentBlue}}p{0.18\textwidth} >{\RaggedRight\arraybackslash}X >{\RaggedRight\arraybackslash}p{0.28\textwidth}}
    \toprule
    \textbf{Tingkat} & \textbf{Aktivitas} & \textbf{Bukti kelulusan tingkat} \\
    \midrule
    1. Diagnostik & Soal pemetaan tanpa bantuan & Baseline per indikator \\
    2. Fondasi & Materi ringkas dan contoh berpikir & Kuis konsep dasar \\
    3. Penerapan & Latihan bertingkat dan variasi konteks & Akurasi topikal \\
    4. Integrasi & Paket subtes dengan batas waktu & Akurasi dan kecepatan \\
    5. Simulasi & Tryout penuh dan evaluasi & Stabilitas hasil serta stamina \\
    \bottomrule
  \end{tabularx}}
  \vspace{0.14cm}
  \centering{\color{StudentNavy}\bfseries\fontsize{10.7}{13.3}\selectfont Siswa yang belum tuntas kembali ke materi dan remedial yang relevan.}
\end{frame}

\begin{frame}{Penalaran Umum dilatih melalui tiga cara berpikir}
  \kicker{Subtes 1 -- Penalaran Umum}
  \begin{columns}[T,onlytextwidth]
    \begin{column}{0.30\textwidth}
      \begin{block}{Induktif}
        Menemukan pola atau aturan umum berdasarkan sejumlah fakta dan hubungan.
      \end{block}
    \end{column}
    \begin{column}{0.30\textwidth}
      \begin{block}{Deduktif}
        Menentukan simpulan yang pasti mengikuti premis atau aturan yang diberikan.
      \end{block}
    \end{column}
    \begin{column}{0.30\textwidth}
      \begin{block}{Kuantitatif}
        Bernalar menggunakan angka, hubungan besaran, dan representasi kuantitatif.
      \end{block}
    \end{column}
  \end{columns}
  \vspace{0.28cm}
  \begin{beamercolorbox}[rounded=true,sep=0.20cm,wd=\textwidth]{block body}
    \centering\bfseries Fokus pembelajaran adalah alasan yang sah, bukan menghafal jenis pola.
  \end{beamercolorbox}
\end{frame}

\begin{frame}{Indikator Penalaran Umum disusun dari sederhana ke kompleks}
  \kicker{Urutan Materi PU}
  {\fontsize{10.1}{12.6}\selectfont
  \renewcommand{\arraystretch}{0.98}
  \begin{tabularx}{\textwidth}{>{\bfseries\color{StudentBlue}}p{0.20\textwidth} >{\RaggedRight\arraybackslash}X >{\RaggedRight\arraybackslash}p{0.30\textwidth}}
    \toprule
    \textbf{Bagian} & \textbf{Materi operasional} & \textbf{Indikator siswa} \\
    \midrule
    Induktif & Pola bilangan, klasifikasi, analogi, kecenderungan, dan generalisasi & Menemukan aturan paling konsisten \\
    Deduktif & Implikasi, negasi, silogisme, kondisi perlu/cukup, dan konsistensi premis & Menentukan simpulan yang valid \\
    Kuantitatif & Perbandingan, estimasi, tabel, grafik, urutan, dan hubungan besaran & Memilih informasi kuantitatif yang relevan \\
    Strategi & Eliminasi opsi, pengujian kasus, dan kontrol waktu & Menjelaskan alasan serta memeriksa simpulan \\
    \bottomrule
  \end{tabularx}}
\end{frame}

\begin{frame}[shrink=20]{PPU melatih pemahaman bahasa melalui konteks}
  \kicker{Subtes 2 -- Pengetahuan dan Pemahaman Umum}
  {\fontsize{10.2}{12.7}\selectfont
  \renewcommand{\arraystretch}{0.99}
  \begin{tabularx}{\textwidth}{>{\bfseries\color{StudentBlue}}p{0.22\textwidth} >{\RaggedRight\arraybackslash}X >{\RaggedRight\arraybackslash}p{0.28\textwidth}}
    \toprule
    \textbf{Domain} & \textbf{Materi} & \textbf{Indikator} \\
    \midrule
    Leksikal & Makna kata, frasa, istilah, rujukan, dan relasi makna & Menentukan makna sesuai konteks \\
    Kalimat & Fungsi unsur, hubungan antarklausa, dan maksud kalimat & Memahami informasi tersurat/tersirat \\
    Paragraf & Gagasan utama, rincian, hubungan, dan simpulan & Memetakan hubungan gagasan \\
    Pengetahuan umum & Informasi sosial, budaya, sains, dan kehidupan akademik & Menggunakan pengetahuan untuk memahami teks \\
    \bottomrule
  \end{tabularx}}
  \vspace{0.12cm}
  \claim{Kosakata dipelajari melalui penggunaan, bukan daftar kata yang terpisah dari bacaan.}
\end{frame}

\begin{frame}{PBM berpusat pada kualitas penyampaian gagasan}
  \kicker{Subtes 3 -- Pemahaman Bacaan dan Menulis}
  \begin{columns}[T,onlytextwidth]
    \begin{column}{0.30\textwidth}
      \begin{block}{Memilih}
        Judul, diksi, kalimat penjelas, serta gagasan yang melengkapi paragraf.
      \end{block}
    \end{column}
    \begin{column}{0.30\textwidth}
      \begin{block}{Menata}
        Urutan kalimat, posisi gagasan, hubungan paragraf, dan penggabungan kalimat.
      \end{block}
    \end{column}
    \begin{column}{0.30\textwidth}
      \begin{block}{Memperbaiki}
        Ejaan, tanda baca, keefektifan, kesejajaran, kerancuan, dan kepaduan.
      \end{block}
    \end{column}
  \end{columns}
  \vspace{0.26cm}
  {\color{StudentBlue}\bfseries\fontsize{11}{13.5}\selectfont Urutan belajar}\par
  {\fontsize{10.8}{13.4}\selectfont Kalimat efektif $\rightarrow$ kepaduan paragraf $\rightarrow$ organisasi teks $\rightarrow$ penyuntingan terpadu.}
\end{frame}

\begin{frame}{Pengetahuan Kuantitatif dimulai dari kelancaran dasar}
  \kicker{Subtes 4 -- Peta Materi PK}
  {\fontsize{10.0}{12.5}\selectfont
  \renewcommand{\arraystretch}{0.98}
  \begin{tabularx}{\textwidth}{>{\bfseries\color{StudentBlue}}p{0.21\textwidth} >{\RaggedRight\arraybackslash}X >{\RaggedRight\arraybackslash}p{0.28\textwidth}}
    \toprule
    \textbf{Domain} & \textbf{Materi inti} & \textbf{Kemampuan} \\
    \midrule
    Bilangan & Operasi, pecahan, persen, rasio, pangkat, akar, dan estimasi & Menghitung serta membandingkan secara efisien \\
    Aljabar & Bentuk aljabar, persamaan, pertidaksamaan, fungsi, dan pola & Memanipulasi hubungan simbolik \\
    Geometri & Bangun datar/ruang, koordinat, sudut, skala, dan pengukuran & Menghubungkan bentuk dan ukuran \\
    Data dan peluang & Tabel, grafik, statistik dasar, pencacahan, dan peluang & Menafsirkan data dan ketidakpastian \\
    \bottomrule
  \end{tabularx}}
\end{frame}

\begin{frame}[shrink=8]{PK menuntut pilihan strategi yang cepat dan tepat}
  \kicker{Indikator Penguasaan}
  \begin{columns}[T,onlytextwidth]
    \begin{column}{0.46\textwidth}
      {\color{StudentBlue}\bfseries Penguasaan konsep}\par
      \begin{itemize}
        \item Mengenali konsep yang relevan.
        \item Menyederhanakan informasi numerik.
        \item Mengubah bentuk representasi.
        \item Memeriksa kewajaran nilai.
      \end{itemize}
    \end{column}
    \begin{column}{0.48\textwidth}
      {\color{StudentBlue}\bfseries Efisiensi prosedur}\par
      \begin{itemize}
        \item Memilih hitung langsung atau estimasi.
        \item Menggunakan perbandingan dan substitusi.
        \item Mengeliminasi pilihan yang tidak mungkin.
        \item Menghindari langkah aljabar yang tidak perlu.
      \end{itemize}
    \end{column}
  \end{columns}
  \vspace{0.18cm}
  \claim{Rumus adalah alat; keputusan memilih alat menjadi bagian dari kemampuan.}
\end{frame}

\begin{frame}{LBI mengembangkan empat proses literasi}
  \kicker{Subtes 5 -- Literasi dalam Bahasa Indonesia}
  \begin{columns}[T,onlytextwidth]
    \begin{column}{0.22\textwidth}
      \begin{block}{Menemukan}
        Mengambil fakta, rincian, rujukan, dan informasi eksplisit.
      \end{block}
    \end{column}
    \begin{column}{0.22\textwidth}
      \begin{block}{Memahami}
        Menghubungkan gagasan, sebab-akibat, tujuan, dan simpulan.
      \end{block}
    \end{column}
    \begin{column}{0.22\textwidth}
      \begin{block}{Menilai}
        Memeriksa bukti, argumen, konsistensi, dan kredibilitas.
      \end{block}
    \end{column}
    \begin{column}{0.22\textwidth}
      \begin{block}{Merefleksikan}
        Menghubungkan isi teks dengan konteks dan persoalan lain.
      \end{block}
    \end{column}
  \end{columns}
  \vspace{0.28cm}
  \centering{\color{StudentNavy}\bfseries\fontsize{14}{17}\selectfont Setiap jawaban harus dapat ditunjukkan kembali pada bukti bacaan.}
\end{frame}

\begin{frame}{Konteks LBI diberi label sejak materi dan latihan}
  \kicker{Saintek, Soshum, dan Personal}
  {\fontsize{10.0}{12.5}\selectfont
  \renewcommand{\arraystretch}{0.98}
  \begin{tabularx}{\textwidth}{>{\bfseries\color{StudentBlue}}p{0.18\textwidth} >{\RaggedRight\arraybackslash}X >{\RaggedRight\arraybackslash}p{0.31\textwidth}}
    \toprule
    \textbf{Konteks} & \textbf{Topik bacaan} & \textbf{Fokus tambahan} \\
    \midrule
    Saintek & Fisika, kimia, biologi, lingkungan, kesehatan, dan teknologi & Istilah, data, proses, dan hubungan sebab-akibat \\
    Soshum & Ekonomi, sejarah, sosiogeografi, masyarakat, dan kebijakan & Argumen, kronologi, kepentingan, dan dampak sosial \\
    Personal & Sastra dan pengalaman manusia & Sudut pandang, gaya, inferensi, makna, dan refleksi \\
    \bottomrule
  \end{tabularx}}
  \vspace{0.14cm}
  \begin{beamercolorbox}[rounded=true,sep=0.17cm,wd=\textwidth]{block body}
    \centering\bfseries Analisis StudentPro memisahkan akurasi LBI saintek dan soshum tanpa mengabaikan bacaan personal.
  \end{beamercolorbox}
\end{frame}

\begin{frame}{Bacaan kebangsaan disiapkan sebagai pengayaan LBI}
  \kicker{Antisipasi 2027 -- Belum Menjadi Ketentuan Resmi}
  \begin{columns}[T,onlytextwidth]
    \begin{column}{0.46\textwidth}
      {\color{StudentBlue}\bfseries Tema yang dapat digunakan}\par
      \begin{itemize}
        \item Konstitusi dan kelembagaan.
        \item Sejarah kebangsaan.
        \item Kebinekaan dan kewargaan.
        \item Kebijakan publik dan masyarakat.
      \end{itemize}
    \end{column}
    \begin{column}{0.48\textwidth}
      {\color{StudentBlue}\bfseries Kemampuan yang tetap dinilai}\par
      \begin{itemize}
        \item Menemukan informasi.
        \item Memahami argumen.
        \item Menilai bukti.
        \item Menyimpulkan dan merefleksikan.
      \end{itemize}
    \end{column}
  \end{columns}
  \vspace{0.16cm}
  \begin{beamercolorbox}[rounded=true,sep=0.18cm,wd=\textwidth]{block body}
    \centering\bfseries Materi tidak menyatakan TWK sebagai subtes baru dan tidak berpusat pada hafalan.
  \end{beamercolorbox}
\end{frame}

\begin{frame}{LBE mengajarkan cara membaca teks akademik}
  \kicker{Subtes 6 -- Literasi dalam Bahasa Inggris}
  {\fontsize{10.2}{12.7}\selectfont
  \renewcommand{\arraystretch}{0.99}
  \begin{tabularx}{\textwidth}{>{\bfseries\color{StudentBlue}}p{0.22\textwidth} >{\RaggedRight\arraybackslash}X >{\RaggedRight\arraybackslash}p{0.28\textwidth}}
    \toprule
    \textbf{Kemampuan} & \textbf{Materi operasional} & \textbf{Indikator} \\
    \midrule
    Locate & Specific information, reference, and textual evidence & Menemukan bukti dengan cepat \\
    Understand & Main idea, detail, relationship, sequence, and inference & Memahami hubungan gagasan \\
    Evaluate & Purpose, tone, argument, evidence, and credibility & Menilai kualitas informasi \\
    Reflect & Implication and connection to another context & Menggunakan isi teks secara bermakna \\
    \bottomrule
  \end{tabularx}}
  \vspace{0.12cm}
  \claim{Grammar dan vocabulary dipelajari sebagai alat memahami teks, bukan tujuan tunggal.}
\end{frame}

\begin{frame}{Urutan LBE menurunkan beban membaca secara bertahap}
  \kicker{Jalur Pembelajaran Bahasa Inggris}
  \begin{columns}[T,onlytextwidth]
    \begin{column}{0.22\textwidth}
      \begin{block}{1. Struktur}
        Mengenali klausa, rujukan, konektor, serta organisasi kalimat.
      \end{block}
    \end{column}
    \begin{column}{0.22\textwidth}
      \begin{block}{2. Paragraf}
        Menentukan gagasan utama, rincian, dan hubungan antarkalimat.
      \end{block}
    \end{column}
    \begin{column}{0.22\textwidth}
      \begin{block}{3. Teks}
        Membaca argumen, inferensi, tujuan, nada, dan bukti.
      \end{block}
    \end{column}
    \begin{column}{0.22\textwidth}
      \begin{block}{4. Waktu}
        Memilih urutan membaca dan menemukan bukti secara efisien.
      \end{block}
    \end{column}
  \end{columns}
  \vspace{0.26cm}
  \centering{\color{StudentNavy}\bfseries\fontsize{13.5}{16.5}\selectfont Bacaan saintek, soshum, dan personal digunakan secara bergantian.}
\end{frame}

\begin{frame}[shrink=2]{Penalaran Matematika menghubungkan masalah dan model}
  \kicker{Subtes 7 -- Tiga Proses Utama}
  \begin{columns}[T,onlytextwidth]
    \begin{column}{0.30\textwidth}
      \begin{block}{Merumuskan}
        Memilih informasi, membuat asumsi, serta mengubah situasi menjadi model matematika.
      \end{block}
    \end{column}
    \begin{column}{0.30\textwidth}
      \begin{block}{Menggunakan}
        Menerapkan konsep, prosedur, representasi, dan strategi penyelesaian.
      \end{block}
    \end{column}
    \begin{column}{0.30\textwidth}
      \begin{block}{Menafsirkan}
        Mengembalikan hasil ke konteks dan memeriksa kewajaran serta keterbatasannya.
      \end{block}
    \end{column}
  \end{columns}
  \vspace{0.27cm}
  \claim{Jawaban matematika belum selesai sebelum maknanya sesuai dengan masalah.}
\end{frame}

\begin{frame}[shrink=10]{Materi PM diorganisasi menurut konten dan konteks}
  \kicker{Peta Penalaran Matematika}
  {\fontsize{10.1}{12.6}\selectfont
  \renewcommand{\arraystretch}{0.98}
  \begin{tabularx}{\textwidth}{>{\bfseries\color{StudentBlue}}p{0.20\textwidth} >{\RaggedRight\arraybackslash}X >{\RaggedRight\arraybackslash}p{0.30\textwidth}}
    \toprule
    \textbf{Dimensi} & \textbf{Cakupan} & \textbf{Fokus latihan} \\
    \midrule
    Konten & Bilangan, aljabar, geometri, pengukuran, data, ketidakpastian, dan perubahan & Pemilihan konsep yang tepat \\
    Pribadi & Keuangan, perjalanan, kesehatan, dan keputusan sehari-hari & Interpretasi bagi individu \\
    Pekerjaan & Produksi, jadwal, biaya, desain, dan efisiensi & Model dan batasan praktis \\
    Sosial/ilmiah & Populasi, lingkungan, kebijakan, eksperimen, dan data & Bukti, tren, dan kewajaran \\
    \bottomrule
  \end{tabularx}}
\end{frame}

\begin{frame}{Strategi waktu menjadi materi lintas subtes}
  \kicker{Kemampuan yang Tidak Berdiri Sendiri}
  \begin{columns}[T,onlytextwidth]
    \begin{column}{0.46\textwidth}
      {\color{StudentBlue}\bfseries Sebelum mengerjakan}\par
      \begin{itemize}
        \item Membaca petunjuk dan mengenali bentuk tugas.
        \item Menetapkan batas waktu internal.
        \item Menentukan urutan membaca atau menghitung.
      \end{itemize}
    \end{column}
    \begin{column}{0.48\textwidth}
      {\color{StudentBlue}\bfseries Saat mengerjakan}\par
      \begin{itemize}
        \item Menandai bukti dan informasi penting.
        \item Melewati hambatan yang tidak proporsional.
        \item Memeriksa kewajaran sebelum berpindah.
      \end{itemize}
    \end{column}
  \end{columns}
  \vspace{0.22cm}
  \begin{beamercolorbox}[rounded=true,sep=0.18cm,wd=\textwidth]{block body}
    \centering\bfseries Kecepatan dilatih setelah cara berpikir benar mulai stabil.
  \end{beamercolorbox}
\end{frame}

\begin{frame}[shrink=16]{Setiap modul StudentPro menggunakan struktur yang sama}
  \kicker{Standar Konten LMS}
  {\fontsize{10.1}{12.6}\selectfont
  \renewcommand{\arraystretch}{0.98}
  \begin{tabularx}{\textwidth}{>{\bfseries\color{StudentBlue}}p{0.20\textwidth} >{\RaggedRight\arraybackslash}X >{\RaggedRight\arraybackslash}p{0.28\textwidth}}
    \toprule
    \textbf{Bagian modul} & \textbf{Isi} & \textbf{Fungsi} \\
    \midrule
    Tujuan & Kemampuan dan indikator yang akan dicapai & Memberi arah belajar \\
    Peta konsep & Hubungan materi dan prasyarat & Menata pengetahuan \\
    Penjelasan & Konsep, strategi, dan kesalahan umum & Membentuk pemahaman \\
    Contoh berpikir & Alasan setiap langkah tanpa membanjiri soal & Memodelkan proses \\
    Cek pemahaman & Kuis singkat dan umpan balik & Menentukan lanjut/remedial \\
    Ringkasan & Pokok materi dan strategi & Memudahkan peninjauan ulang \\
    \bottomrule
  \end{tabularx}}
\end{frame}

\begin{frame}{Penandaan konten membuat analisis lebih bermakna}
  \kicker{Metadata Materi dan Soal}
  \begin{columns}[T,onlytextwidth]
    \begin{column}{0.46\textwidth}
      \begin{itemize}
        \item Subtes dan bagian kemampuan.
        \item Domain materi dan indikator.
        \item Tingkat penguasaan.
        \item Tingkat kesulitan.
      \end{itemize}
    \end{column}
    \begin{column}{0.48\textwidth}
      \begin{itemize}
        \item Konteks bacaan atau masalah.
        \item Jenis proses berpikir.
        \item Estimasi waktu.
        \item Jenis kesalahan yang mungkin.
      \end{itemize}
    \end{column}
  \end{columns}
  \vspace{0.22cm}
  \claim{Satu nilai dapat diuraikan menjadi materi, konteks, proses, waktu, dan pola kesalahan.}
\end{frame}

\begin{frame}[shrink=12]{Ketuntasan materi selalu diikuti keputusan}
  \kicker{Asesmen dan Tindak Lanjut}
  {\fontsize{10.0}{12.5}\selectfont
  \renewcommand{\arraystretch}{0.98}
  \begin{tabularx}{\textwidth}{>{\bfseries\color{StudentBlue}}p{0.20\textwidth} >{\RaggedRight\arraybackslash}X >{\RaggedRight\arraybackslash}p{0.31\textwidth}}
    \toprule
    \textbf{Hasil} & \textbf{Interpretasi} & \textbf{Keputusan} \\
    \midrule
    Belum memahami & Konsep atau prasyarat belum kuat & Ulang materi dan bantuan terarah \\
    Paham tetapi lambat & Proses benar, belum otomatis/efisien & Latihan bertahap dengan batas waktu \\
    Cepat tetapi tidak teliti & Strategi waktu terlalu agresif & Latihan verifikasi dan kontrol keputusan \\
    Stabil dan tuntas & Akurasi serta waktu memenuhi target internal & Pengayaan dan paket integrasi \\
    Turun pada tryout penuh & Stamina atau perpindahan subtes bermasalah & Simulasi bertahap dan evaluasi ritme \\
    \bottomrule
  \end{tabularx}}
\end{frame}

\begin{frame}[shrink=6]{Orang tua menerima peta kemajuan yang sederhana}
  \kicker{Komunikasi Kurikulum}
  \begin{columns}[T,onlytextwidth]
    \begin{column}{0.46\textwidth}
      {\color{StudentBlue}\bfseries Yang dilaporkan}\par
      \begin{itemize}
        \item Materi yang sudah dipelajari.
        \item Indikator yang sudah atau belum tuntas.
        \item Tren akurasi dan waktu.
        \item Remedial atau pengayaan berikutnya.
      \end{itemize}
    \end{column}
    \begin{column}{0.48\textwidth}
      {\color{StudentBlue}\bfseries Yang tidak dilakukan}\par
      \begin{itemize}
        \item Membandingkan anak dengan siswa lain.
        \item Menarik kesimpulan dari satu nilai.
        \item Menganggap semua kelemahan sebagai kurang belajar.
        \item Mengubah target berdasarkan rumor format.
      \end{itemize}
    \end{column}
  \end{columns}
  \vspace{0.16cm}
  \begin{beamercolorbox}[rounded=true,sep=0.18cm,wd=\textwidth]{block body}
    \centering\bfseries Orang tua melihat arah perbaikan, bukan hanya daftar nilai.
  \end{beamercolorbox}
\end{frame}

\begin{frame}{Kesimpulan peta materi UTBK}
  \kicker{Arah Kurikulum StudentPro}
  \begin{columns}[T,onlytextwidth]
    \begin{column}{0.30\textwidth}
      \begin{block}{Lengkap}
        Tujuh subtes dipetakan menjadi domain, indikator, dan proses berpikir.
      \end{block}
    \end{column}
    \begin{column}{0.30\textwidth}
      \begin{block}{Bertahap}
        Diagnostik, fondasi, penerapan, integrasi, dan simulasi.
      \end{block}
    \end{column}
    \begin{column}{0.30\textwidth}
      \begin{block}{Adaptif}
        Urutan dan tindak lanjut mengikuti data penguasaan setiap siswa.
      \end{block}
    \end{column}
  \end{columns}
  \vspace{0.28cm}
  \centering{\color{StudentNavy}\bfseries\fontsize{14}{17}\selectfont Kurikulum yang sama menyediakan jalur belajar yang berbeda bagi setiap siswa.}
\end{frame}

\begin{frame}{Sumber resmi materi}
  \kicker{Rujukan Bab 6 -- Bagian 1}
  {\color{StudentBlue}\bfseries\fontsize{14.5}{17.5}\selectfont 1. Panitia SNPMB}\par
  \vspace{0.05cm}
  {\fontsize{11}{13.5}\selectfont Framework UTBK--SNBT Tahun 2026.}\par
  {\fontsize{10.5}{13}\selectfont\href{https://snpmb.id/fr/}{https://snpmb.id/fr/}}\par
  {\color{StudentGray}\fontsize{10}{12.4}\selectfont Sumber tujuan, materi, proses, konteks, jumlah soal, dan durasi tujuh subtes.}\par
  \vspace{0.34cm}
  {\color{StudentBlue}\bfseries\fontsize{14.5}{17.5}\selectfont 2. Panitia SNPMB}\par
  \vspace{0.05cm}
  {\fontsize{11}{13.5}\selectfont Informasi Umum UTBK--SNBT 2026.}\par
  {\fontsize{10.5}{13}\selectfont\href{https://snpmb.id/utbk-snbt/informasi-umum}{snpmb.id/utbk-snbt/informasi-umum}}\par
  {\color{StudentGray}\fontsize{10}{12.4}\selectfont Digunakan untuk menjaga batas antara acuan resmi 2026 dan persiapan 2027.}\par
\end{frame}

\begin{frame}{Sumber internal dan catatan status}
  \kicker{Rujukan Bab 6 -- Bagian 2}
  {\color{StudentBlue}\bfseries\fontsize{14.5}{17.5}\selectfont 3. Bank materi dan soal StudentPro}\par
  \vspace{0.05cm}
  {\fontsize{10.8}{13.2}\selectfont Struktur modul, indikator, kuis topikal, paket subtes, analisis, remedial, dan pengayaan.}\par
  {\color{StudentGray}\fontsize{9.8}{12}\selectfont Digunakan sebagai dasar implementasi kurikulum di LMS.}\par
  \vspace{0.28cm}
  {\color{StudentBlue}\bfseries\fontsize{14.5}{17.5}\selectfont 4. Antisipasi bacaan kebangsaan dalam LBI}\par
  \vspace{0.05cm}
  {\fontsize{10.8}{13.2}\selectfont Pengayaan konteks literasi internal; bukan penetapan TWK sebagai subtes UTBK 2027.}\par
  \vfill
  \begin{beamercolorbox}[rounded=true,sep=0.16cm,wd=\textwidth]{block body}
    \centering\bfseries Bab 7 menerapkan peta materi ini dalam contoh soal dan pembahasan.
  \end{beamercolorbox}
\end{frame}
